\subsection{Metric definitions and identities}
All recognition claims use two threshold-free metrics.
\begin{definition}[AUROC]
For scores $s_i$ and labels $y_i\in\{0,1\}$,
$\mathrm{AUROC} = P(s_i > s_j \mid y_i{=}1, y_j{=}0) + \tfrac12 P(s_i = s_j)$,
the area under the receiver operating characteristic curve traced by varying
the decision threshold.
\end{definition}
\begin{proposition}[Mann--Whitney identity]\label{prop:mw}
Let $U$ be the number of positive--negative pairs with $s_i > s_j$ plus half
the ties. Then $\mathrm{AUROC} = U/(n_+ n_-)$ exactly, and under the null of
no label information $\mathrm{AUROC}$ concentrates around $1/2$ with
variance $(n_+ + n_- + 1)/(12 n_+ n_-)$.
\end{proposition}
\begin{proof}[Proof sketch]
The ROC curve is the empirical survival function of positive scores against
that of negatives; integrating by parts turns the area into the pairwise
probability. The variance formula is the standard Mann--Whitney null variance
under exchangeability of the pooled sample.
\end{proof}
\begin{definition}[AUPRC]
With precision $\pi(t) = P(y{=}1 \mid s \ge t)$ and recall
$\rho(t) = P(s \ge t \mid y{=}1)$, $\mathrm{AUPRC} = \int_0^1 \pi\, d\rho$,
estimated by average precision.
\end{definition}
AUPRC is the more sensitive metric under class imbalance because it ignores
true negatives entirely; our test split is balanced, so the two metrics track
each other, and the observed AUPRC gaps (e.g.\ GNN $0.680$ vs CNN $0.879$ on
the scale-up test) are the more informative signal of ranking quality in the
high-score region that the design and screen stages consume.

\subsection{Operating points}
The 5\% FPR operating point used for APD6 sensitivity is chosen on the
held-out negatives only: the threshold is the 95th percentile of ensemble
scores over the 2{,}925 test negatives, then applied to the APD6 positives.
Because the threshold never touches the positive scores, the sensitivity
estimate is an honest one-sided generalization measurement rather than a
fitted operating point.
